\chapter{Pierwszy dokument}
\label{cha:pierwszyDokument}

W rozdziale tym przedstawiono podstawowe informacje dotyczące struktury prostych plików \LaTeX a. Omówiono również metody kompilacji plików z zastosowaniem programów \emph{latex} oraz \emph{pdflatex}.

%---------------------------------------------------------------------------

\section{Struktura dokumentu}
\label{sec:strukturaDokumentu}

Plik \LaTeX owy jest plikiem tekstowym, który oprócz tekstu zawiera polecenia formatujące ten tekst (analogicznie do języka HTML). Plik składa się z dwóch części:
\begin{enumerate}%[1)]
      \item Preambuły -- określającej klasę dokumentu oraz zawierającej m.in. polecenia dołączającej dodatkowe pakiety;

      \item Części głównej -- zawierającej zasadniczą treść dokumentu.
\end{enumerate}

% ------------------------
\begin{lstlisting}
\documentclass[a4paper,12pt]{article}      % preambuła
\usepackage[polish]{babel}
\usepackage[utf8]{inputenc}
\usepackage[T1]{fontenc}
\usepackage{times}

\begin{document}                           % część główna

\section{Sztuczne życie}


\end{document}
\end{lstlisting}

Nie ma żadnych przeciwskazań do tworzenia dokumentów w~\LaTeX u w~języku polskim. Plik źródłowy jest zwykłym plikiem tekstowym i~do jego przygotowania można użyć dowolnego edytora tekstów, a~polskie znaki wprowadzać używając prawego klawisza \texttt{Alt}. Jeżeli po kompilacji dokumentu polskie znaki nie są wyświetlane poprawnie, to na 95\% źle określono sposób kodowania znaków (należy zmienić opcje wykorzystywanych pakietów).
%---------------------------------------------------------------------------


% ------------------------
\section{Kompilacja}
\label{sec:kompilacja}


Załóżmy, że przygotowany przez nas dokument zapisany jest w pliku \texttt{test.tex}. Kolejno wykonane poniższe polecenia (pod warunkiem, że w pierwszym przypadku nie wykryto błędów i kompilacja zakończyła się sukcesem) pozwalają uzyskać nasz dokument w formacie pdf:
\begin{lstlisting}
latex test.tex
dvips test.dvi -o test.ps
ps2pdf test.ps
\end{lstlisting}
%
lub za pomocą PDF\LaTeX:
\begin{lstlisting}
pdflatex test.tex
\end{lstlisting}

Przy pierwszej kompilacji po zmiane tekstu, dodaniu nowych etykiet itp., \LaTeX~tworzy sobie spis rozdziałów, obrazków, tabel itp., a dopiero przy następnej kompilacji korzysta z tych informacji.

W pierwszym przypadku rysunki powinny być przygotowane w~formacie eps, a~w~drugim w~formacie pdf. Ponadto, jeżeli używamy polecenia \texttt{pdflatex test.tex} można wstawiać grafikę bitową (np. w formacie jpg).



%---------------------------------------------------------------------------
\section{Narzędzia}
\label{sec:narzedzia}

Do tworzenia dokumentów w systemie \LaTeX{} dostępnych jest wiele narzędzi, które wspomagają przygotowanie, edycję, kompilację oraz zarządzanie dokumentami. Wybór odpowiedniego środowiska zależy przede wszystkim od preferowanego sposobu pracy, wymagań projektu oraz doświadczenia użytkownika.

Narzędzia do pracy z systemem \LaTeX{} można podzielić na trzy podstawowe grupy: edytory online, edytory lokalne (offline) oraz programy wspomagające zarządzanie bibliografią. W przypadku pracy lokalnej konieczne jest również zainstalowanie odpowiedniej dystrybucji systemu \LaTeX{}.

\subsection{Edytory online}

Edytory online umożliwiają przygotowywanie i kompilowanie dokumentów bez konieczności instalowania dodatkowego oprogramowania na komputerze. Do pracy wystarczy przeglądarka internetowa oraz dostęp do sieci.

\begin{itemize}

    \item \textbf{Overleaf} -- jedno z najpopularniejszych środowisk do tworzenia dokumentów w systemie \LaTeX{}. Umożliwia edycję i kompilację dokumentów bezpośrednio w przeglądarce internetowej, a także współpracę kilku użytkowników nad jednym projektem. Najważniejsze funkcje:
    \begin{itemize}
        \item automatyczna kompilacja dokumentów i podgląd plików PDF,
        \item obsługa kompilatorów pdfLaTeX, XeLaTeX oraz LuaLaTeX,
        \item obsługa bibliografii z wykorzystaniem BibTeX i Biber,
        \item współpraca nad dokumentem w czasie rzeczywistym,
        \item dostęp do gotowych szablonów prac dyplomowych, artykułów naukowych, prezentacji i innych dokumentów,
        \item możliwość importowania i eksportowania projektów,
        \item integracja z systemami kontroli wersji Git i GitHub (w ramach odpowiednich planów płatnych).
    \end{itemize}

    Overleaf oferuje plan bezpłatny, który jest wystarczający do przygotowania większości podstawowych dokumentów. Bardziej zaawansowane funkcje, w tym rozszerzona współpraca i integracja z GitHub, wymagają odpowiedniej subskrypcji.

    Strona internetowa: \url{https://www.overleaf.com/} \cite{Overleaf}.


\item \textbf{OpenAI Prism} -- internetowe środowisko do tworzenia dokumentów w systemie \LaTeX{}, opracowane przez OpenAI. Łączy tradycyjną edycję i kompilację dokumentów z funkcjami sztucznej inteligencji, wspomagającymi przygotowywanie tekstów naukowych, artykułów i prac dyplomowych.

Wśród najważniejszych funkcji można wyróżnić:
\begin{itemize}
    \item edycję i kompilację dokumentów \LaTeX{} online,
    \item współpracę wielu użytkowników w czasie rzeczywistym,
    \item generowanie i poprawianie kodu \LaTeX{} z pomocą AI,
    \item wykrywanie i analizowanie błędów kompilacji,
    \item wspomaganie redagowania i formatowania tekstów naukowych,
    \item wyszukiwanie literatury i zarządzanie cytowaniami,
    \item integrację z narzędziem Zotero.
\end{itemize}

Prism oferuje bezpłatny dostęp do podstawowych funkcji za pośrednictwem konta ChatGPT. Nie wymaga instalowania lokalnej dystrybucji \LaTeX{}. 

Strona internetowa: \url{https://prism.openai.com/} \cite{PrismExample}.

    \item \textbf{Papeeria} -- internetowe środowisko do tworzenia dokumentów w systemie \LaTeX{}, umożliwiające pracę bez instalowania dodatkowego oprogramowania.

    Najważniejsze funkcje:
    \begin{itemize}
        \item edycja dokumentów bezpośrednio w przeglądarce,
        \item automatyczna kompilacja dokumentów,
        \item współpraca użytkowników nad projektem,
        \item podpowiadanie składni i poleceń,
        \item możliwość korzystania z gotowych szablonów,
        \item obsługa dokumentów wykorzystujących LaTeX i Markdown.
    \end{itemize}

    Papeeria stanowi alternatywę dla Overleaf i może być wykorzystywana do przygotowywania prostych oraz bardziej rozbudowanych dokumentów.

    Strona internetowa: \url{https://www.papeeria.com/} \cite{PapeeriaExample}.

\end{itemize}


\subsection{Edytory lokalne (offline)}

Edytory lokalne umożliwiają tworzenie i edytowanie dokumentów bez konieczności stałego dostępu do Internetu. W przeciwieństwie do środowisk online wymagają zainstalowania odpowiedniego oprogramowania oraz dystrybucji systemu \LaTeX{}, która odpowiada za kompilację dokumentów.

\begin{itemize}

    \item \textbf{Visual Studio Code (VS Code)} -- uniwersalny edytor kodu źródłowego, który dzięki rozszerzeniu \textit{LaTeX Workshop} może pełnić funkcję rozbudowanego środowiska do tworzenia dokumentów w systemie \LaTeX{}. Najważniejsze funkcje:
    \begin{itemize}
        \item podświetlanie składni i automatyczne uzupełnianie poleceń,
        \item kompilacja dokumentów i podgląd plików PDF,
        \item synchronizacja pomiędzy kodem źródłowym a dokumentem PDF,
        \item wykrywanie błędów składniowych,
        \item obsługa bibliografii BibTeX i Biber,
        \item integracja z systemem kontroli wersji Git,
        \item obsługa rozbudowanych projektów składających się z wielu plików.
    \end{itemize}

    Aby rozpocząć pracę z dokumentami \LaTeX{} w Visual Studio Code, należy:
    \begin{enumerate}
        \item zainstalować Visual Studio Code,
        \item zainstalować dystrybucję \LaTeX{}, np. MiKTeX lub TeX Live,
        \item zainstalować rozszerzenie LaTeX Workshop,
        \item skonfigurować środowisko kompilacji, jeśli ustawienia domyślne nie działają poprawnie.
    \end{enumerate}

    Visual Studio Code jest szczególnie polecany osobom, które korzystają już z tego środowiska podczas programowania lub chcą zarządzać dokumentami z wykorzystaniem systemu Git.

    Strona internetowa: \url{https://code.visualstudio.com/}

    Rozszerzenie LaTeX Workshop: \url{https://github.com/James-Yu/LaTeX-Workshop}

    Instrukcja instalacji i konfiguracji rozszerzenia LaTeX Workshop oraz kompilacji dokumentów PDF w środowisku Visual Studio Code (Windows) jest dostępna w materiale wideo: \url{https://www.youtube.com/watch?v=4lyHIQl4VM8} \cite{LatexYT}.


    \item \textbf{TeXstudio} -- bezpłatne, wieloplatformowe środowisko przeznaczone bezpośrednio do tworzenia dokumentów w systemie \LaTeX{}.

    Najważniejsze funkcje:
    \begin{itemize}
        \item podświetlanie składni i automatyczne uzupełnianie poleceń,
        \item zintegrowany podgląd dokumentów PDF,
        \item synchronizacja kodu źródłowego z dokumentem PDF,
        \item sprawdzanie pisowni, odwołań i składni,
        \item kreatory tabel, równań matematycznych oraz wstawiania grafik,
        \item obsługa bibliografii i odwołań,
        \item możliwość konfigurowania procesu kompilacji dokumentów.
    \end{itemize}

    TeXstudio jest dobrym rozwiązaniem zarówno dla początkujących, jak i zaawansowanych użytkowników. W przeciwieństwie do Visual Studio Code zostało zaprojektowane specjalnie do pracy z dokumentami \LaTeX{}, dzięki czemu wiele potrzebnych funkcji jest dostępnych bez instalowania dodatkowych rozszerzeń.

    Strona internetowa: \url{https://www.texstudio.org/} \cite{TeXstudioEditor}.

\end{itemize}


\subsection{Dystrybucje systemu LaTeX}

W przypadku korzystania z edytorów lokalnych niezbędna jest instalacja odpowiedniej dystrybucji systemu \LaTeX{}. Dystrybucja zawiera kompilatory, pakiety, czcionki oraz narzędzia potrzebne do przekształcania plików źródłowych \texttt{.tex} w gotowe dokumenty, najczęściej w formacie PDF.

Do najpopularniejszych dystrybucji należą:

\begin{itemize}

    \item \textbf{MiKTeX} -- dystrybucja dostępna m.in. dla systemu Windows. Umożliwia instalowanie brakujących pakietów na żądanie, co ułatwia rozpoczęcie pracy z dokumentami \LaTeX{}.

    Strona internetowa: \url{https://miktex.org/} \cite{MiKTeX}.

    \item \textbf{TeX Live} -- rozbudowana, wieloplatformowa dystrybucja systemu \LaTeX{}, dostępna dla systemów Windows, Linux i macOS. Zawiera szeroki zestaw pakietów oraz narzędzi do przygotowywania dokumentów.

    Strona internetowa: \url{https://www.tug.org/texlive/} \cite{Tex}.

    \item \textbf{MacTeX} -- dystrybucja przeznaczona dla systemu macOS, oparta na TeX Live. Zawiera narzędzia potrzebne do kompilacji dokumentów oraz dodatkowe aplikacje wspomagające pracę z systemem \LaTeX{}.

    Strona internetowa: \url{https://www.tug.org/mactex/} \cite{MacTeX}.

\end{itemize}

Należy pamiętać, że edytor dokumentów i dystrybucja systemu \LaTeX{} pełnią różne funkcje. Edytor służy przede wszystkim do przygotowywania kodu źródłowego dokumentu, natomiast dystrybucja dostarcza narzędzi niezbędnych do jego kompilacji.


\subsection{Narzędzia do zarządzania bibliografią}

Podczas przygotowywania prac dyplomowych i publikacji naukowych istotnym elementem jest prawidłowe zarządzanie źródłami bibliograficznymi. W systemie \LaTeX{} bibliografia jest często przechowywana w plikach \texttt{.bib}, wykorzystywanych przez BibTeX lub BibLaTeX wraz z programem Biber.

Do popularnych narzędzi wspomagających zarządzanie bibliografią należą:

\begin{itemize}

    \item \textbf{Zotero} -- bezpłatne narzędzie do gromadzenia, organizowania i zarządzania literaturą naukową.

    Najważniejsze funkcje:
    \begin{itemize}
        \item importowanie danych bibliograficznych z baz naukowych i stron internetowych,
        \item organizowanie publikacji w kolekcje,
        \item przechowywanie i oznaczanie dokumentów PDF,
        \item eksport bibliografii do formatów BibTeX i BibLaTeX,
        \item synchronizacja bibliotek pomiędzy urządzeniami,
        \item możliwość współdzielenia zbiorów bibliograficznych.
    \end{itemize}

    Integrację Zotero z projektami \LaTeX{} można rozszerzyć za pomocą dodatku \textit{Better BibTeX}, który umożliwia m.in. automatyczną aktualizację eksportowanych plików bibliograficznych.

    Strona internetowa: \url{https://www.zotero.org/}

    Dodatek Better BibTeX:
    \url{https://retorque.re/zotero-better-bibtex/}

    \item \textbf{JabRef} -- bezpłatne narzędzie typu open source, zaprojektowane z myślą o zarządzaniu bibliografią w formatach BibTeX i BibLaTeX.

    Najważniejsze funkcje:
    \begin{itemize}
        \item tworzenie i edycja plików \texttt{.bib},
        \item organizowanie i wyszukiwanie pozycji bibliograficznych,
        \item importowanie metadanych publikacji,
        \item zarządzanie kluczami cytowań,
        \item wykrywanie duplikatów i uzupełnianie danych bibliograficznych,
        \item bezpośrednia współpraca z projektami \LaTeX{}.
    \end{itemize}

    JabRef jest szczególnie przydatny podczas przygotowywania większych dokumentów naukowych, w których bibliografia jest zarządzana bezpośrednio za pomocą plików \texttt{.bib}.

    Strona internetowa: \url{https://www.jabref.org/}

    \item \textbf{Mendeley Reference Manager} -- narzędzie do organizowania publikacji naukowych, przechowywania dokumentów PDF oraz zarządzania cytowaniami.

    Najważniejsze funkcje:
    \begin{itemize}
        \item gromadzenie i katalogowanie publikacji naukowych,
        \item importowanie metadanych bibliograficznych,
        \item organizowanie dokumentów PDF,
        \item dodawanie notatek i oznaczeń w publikacjach,
        \item eksport danych bibliograficznych do formatu BibTeX,
        \item synchronizacja biblioteki użytkownika.
    \end{itemize}

    Współpraca Mendeley Reference Manager z systemem \LaTeX{} jest możliwa poprzez eksport danych bibliograficznych do pliku \texttt{.bib}. Eksport ten jest dostępny w aplikacji komputerowej.

    Strona internetowa: \url{https://www.mendeley.com/}

\end{itemize}


\subsection{Wybór odpowiedniego środowiska}

Wybór narzędzi do pracy z systemem \LaTeX{} powinien uwzględniać charakter przygotowywanego dokumentu, doświadczenie użytkownika oraz preferowany sposób pracy.

Dla osób rozpoczynających naukę \LaTeX{} rekomendowanym rozwiązaniem jest \textbf{Overleaf}, który nie wymaga instalowania ani konfigurowania lokalnego środowiska. Umożliwia szybkie rozpoczęcie pracy oraz korzystanie z gotowych szablonów dokumentów.

Osobom preferującym pracę lokalną można polecić \textbf{TeXstudio}, które oferuje zintegrowane środowisko edycji i kompilacji dokumentów. Alternatywnym rozwiązaniem jest \textbf{Visual Studio Code z rozszerzeniem LaTeX Workshop}, szczególnie przydatny w przypadku projektów wykorzystujących system kontroli wersji Git.

Do zarządzania bibliografią polecane są przede wszystkim \textbf{Zotero z dodatkiem Better BibTeX} oraz \textbf{JabRef}. Oba rozwiązania umożliwiają wygodne zarządzanie źródłami wykorzystywanymi w pracach dyplomowych i publikacjach naukowych.

Niezależnie od wybranego środowiska należy pamiętać o regularnym zapisywaniu kopii zapasowych projektu, zachowaniu uporządkowanej struktury plików oraz weryfikowaniu poprawności danych bibliograficznych.

%---------------------------------------------------------------------------

\section{Przygotowanie dokumentu}
\label{sec:przygotowanieDokumentu}

Dokumenty w systemie \LaTeX{} tworzone są na podstawie plików źródłowych z rozszerzeniem \texttt{.tex}, które zawierają zwykły tekst oraz polecenia odpowiedzialne za jego strukturę i formatowanie. W przeciwieństwie do tradycyjnych edytorów tekstu, takich jak Microsoft Word, w systemie \LaTeX{} wygląd dokumentu jest generowany  podczas kompilacji pliku źródłowego.

Podczas przygotowywania dokumentu należy pamiętać o kilku podstawowych zasadach:

\begin{itemize}

    \item \textbf{Spacje pomiędzy wyrazami} -- w kodzie źródłowym można stosować pojedyncze lub wielokrotne spacje. Podczas kompilacji kolejne spacje są zazwyczaj interpretowane jako pojedynczy odstęp.

    Aby zapobiec rozdzieleniu dwóch elementów tekstu pomiędzy wierszami, należy zastosować tzw. \textit{twardą spację}, zapisywaną za pomocą znaku tyldy (\texttt{\textasciitilde}). Rozwiązanie to jest szczególnie przydatne np. przy zapisie odwołań do rysunków, tabel i rozdziałów.

    \item \textbf{Tworzenie akapitów} -- nowy akapit rozpoczyna się poprzez pozostawienie pustej linii w kodzie źródłowym. Pojedyncze przejście do nowej linii nie powoduje utworzenia nowego akapitu. Liczba kolejnych pustych linii nie ma znaczenia dla podziału tekstu na akapity.

    \item \textbf{Automatyczne formatowanie tekstu} -- system \LaTeX{} automatycznie odpowiada za skład dokumentu,
    w tym rozmieszczenie tekstu, podział na wiersze, wyrównanie akapitów oraz numerację elementów.

    Nie należy ręcznie poprawiać układu tekstu za pomocą dodatkowych spacji lub pustych linii. Wszelkie zmiany formatowania powinny być wprowadzane za pomocą odpowiednich poleceń i mechanizmów systemu \LaTeX{}.

\end{itemize}

\noindent
\textbf{Przydatne materiały i narzędzia wspomagające pracę z systemem \LaTeX{}:}

\begin{itemize}

    \item \textbf{Overleaf Documentation} --
    dokumentacja oraz materiały edukacyjne dotyczące składni,
    formatowania dokumentów, wstawiania grafik, tworzenia tabel
    i zarządzania bibliografią.
    Dostęp: \url{https://www.overleaf.com/learn}.

    \item \textbf{CodeCogs Equation Editor} --
    internetowy edytor równań matematycznych, umożliwiający
    tworzenie i podgląd wzorów oraz generowanie kodu \LaTeX{},
    który można wykorzystać w dokumencie.
    Dostęp: \url{https://editor.codecogs.com/}.

    \item \textbf{Tables Generator} --
    internetowe narzędzie wspomagające projektowanie tabel.
    Umożliwia przygotowanie struktury tabeli i wygenerowanie
    odpowiadającego jej kodu \LaTeX{}, który następnie można
    umieścić w dokumencie.
    Dostęp: \url{https://www.tablesgenerator.com/}.

\end{itemize}

Wymienione narzędzia mogą ułatwić przygotowywanie dokumentów,
szczególnie osobom rozpoczynającym naukę systemu \LaTeX{}.
Należy jednak pamiętać, że wygenerowany kod powinien zostać
sprawdzony i w razie potrzeby dostosowany do struktury
oraz wymagań formatowania przygotowywanego dokumentu.
